\documentclass[11pt]{article}
\usepackage[T1]{fontenc}
\usepackage{lmodern,amsmath,amssymb,amsthm,mathtools,booktabs}
\usepackage[margin=1in]{geometry}
\usepackage[colorlinks=true,linkcolor=blue!50!black,urlcolor=blue!50!black,citecolor=blue!50!black]{hyperref}
\usepackage{xcolor}
\usepackage{microtype}
\newtheorem{theorem}{Theorem}[section]
\newtheorem{lemma}[theorem]{Lemma}
\newtheorem{corollary}[theorem]{Corollary}
\newtheorem{proposition}[theorem]{Proposition}
\theoremstyle{remark}\newtheorem{remark}[theorem]{Remark}
\newcommand{\U}{\mathcal U}
\newcommand{\E}{\mathbb E}
\newcommand{\Pp}{\mathbb P}
\newcommand{\abs}[1]{\left|#1\right|}
\newcommand{\hG}{\widehat G}
\newcommand{\id}{\mathrm{id}}
\newcommand{\GCD}{\operatorname{gcd}}
\newcommand{\LCM}{\operatorname{lcm}}
\newcommand{\supp}{\operatorname{supp}}
\title{Small-prime factorizations at arbitrary retained density\\\large A growing-rank quotient law from Dirichlet GRH}
\author{Research note for the ProveIt project}
\date{27 November 2026}
\hypersetup{pdftitle={Small-prime factorizations at arbitrary retained density},pdfauthor={Research note for ProveIt},pdfsubject={Dirichlet GRH and growing-label quotient laws}}
\begin{document}
\maketitle
\begin{abstract}
Assume Dirichlet GRH, in the precise form exported by the vendored \texttt{grh-lean} development. Fix a retained density $a>0$ and accuracy exponent $A>0$. We construct a single list of $O_{a,A}(R\log Q)$ distinct primes of size $O_{a,A}(\log^2 Q)$, valid for all choices of at most $R$ moduli up to $Q$, where $R\ll_{a,A}\log Q/\log\log Q$. After arbitrary deletions leaving density at least $a$, the ordered factorizations with prescribed block sizes $\gg_{a,A}\log Q$ obey a relative $Q^{-A}$ equidistribution law on precisely the occupied fibres of a finite character quotient. The extra quotient has order at most $\lfloor1/a\rfloor^{r-1}$ beyond the compulsory total-product congruence. This index bound is sharp. The principal new ingredient is a cycle-compression lemma for integer transportation tables: each attainable quotient value can be reached from any assignment by moving $O_a(r)$ primes, rather than exponentially many in $r$. This gives polynomially large occupied quotient masses and an exact support theorem at the growing-label scale. The result strengthens the fixed-label minority-density theorem in the merged ProveIt article, without requiring GRH for any non-Dirichlet $L$-function.
\end{abstract}
\section{Input, scope and comparison}
The vendored readable Lean package \cite{vendor,statements,main} is reported to be byte-identical to upstream commit
\begin{center}\small\texttt{72fead102f663d5c7a407964a13dc0cf9d977207}.\end{center}
The Mathlib-facing definitions in \texttt{src/Statements.lean} and theorem bodies in \texttt{src/Main.lean} state, in particular, nonvanishing for $\Re s>1/2$ (except the principal pole) and critical-line location of zeros in $0<\Re s<1$, for every Dirichlet character of every positive modulus, including imprimitive characters. The names are \texttt{LFunction\_ne\_zero\_of\_half\_lt\_re} and \texttt{LFunction\_eq\_zero\_re}. The README reports Lean 4.35.0, Mathlib revision \texttt{56a005855ea17e0b0958f73ec801ae140b0a1679}, only the usual three axioms, and independent checker acceptance. This note reads the sources but does not repeat the build or independent kernel checks. The result below is proved as a mathematical implication from this Dirichlet GRH input. It can be treated as unconditional if that input theorem is accepted. It is not a new proof of GRH, nor a claim about nonabelian Artin or general automorphic functions.

\paragraph{Mechanism.} The files \texttt{New/Reduction/MainInequality.lean} and \texttt{OneSidedBound.lean} \cite{ineq,oneside} compare two bounds on one corrected squared covolume; the central inequality is $m\,\mathrm{rem}\le fa(\log b)^5$ in the source notation. The remainder is obtained from a nonnegative, smoothed prime-power statistic. With a large second cutoff the proof gets a one-sided estimate for
\[
R_{d,q}(x)=\sum_{n\le x}\Lambda(n)\mathbf1_{n\equiv1\bmod q}(1-\chi_d(n))
\left(\frac1n-\frac1x\right)-\frac{\log x}{\varphi(q)}-c_{d,q}.
\]
Its analytic criterion excludes off-line zeros after choosing an auxiliary twist. Our application uses the theorem exported by this chain, not its fixed-modulus onset estimates as though they were uniform in a growing conductor. The uniform conductor dependence needed here is supplied by the classical smoothed explicit formula under Dirichlet GRH.

\paragraph{What is new relative to the article.} The merged article \cite{article} contains 131 earlier reports; its introduction and source section describe the scope of the input and its analytic mechanism. We inspected the relevant theorem headings in its analytic, zero, recovery, spanning and direct-arithmetic sections, and the full relevant factorization subsection of \texttt{sections/dir.tex} \cite{dir}. In particular:
\begin{center}
\small
\begin{tabular}{@{}p{0.29\linewidth}p{0.29\linewidth}p{0.33\linewidth}@{}}
\toprule
Feature & Existing simultaneous theorem & Existing minority theorem / this note\\
\midrule
Retained density & $>1/2$ by a fixed margin & Any fixed $a>0$\\
Number of factors & $r\ll\log Q/\log\log Q$ & Fixed $r$ / growing $r\ll\log Q/\log\log Q$\\
Moduli & All $q_j\le Q$ & All $q_j\le Q$\\
Extra obstruction & None & Quotient index $\le\lfloor1/a\rfloor^{r-1}$\\
Prime height & $O(\log^2 Q)$ & Same height, uniformly in growing $r$\\
\bottomrule
\end{tabular}
\end{center}
The earlier minority theorem (label \texttt{dir:thm:minority}) uses small bias up to $Q^r$ and a correction argument requiring up to the order of the whole obstruction group in moves. Neither bound gives the desired growing-$r$ result. We replace the first by bounded-word comparisons up to $Q^{4\lfloor1/a\rfloor}$, and the second by a transportation-cycle argument with only $O_a(r)$ moves. These are the two essential changes. The multi-label Fourier estimate itself is adapted from the article. Published work on subset-product covering includes Booker--Pomerance \cite{bp}, Klurman--Shparlinski--Ter\"av\"ainen \cite{kst} and Hughes \cite{hughes}; those results concern different coverage/size regimes. Our targeted literature search did not establish priority beyond this comparison. The theorem is a candidate novel contribution, not a certified claim of first discovery.

\section{The growing-label theorem}
For an integer $q\ge1$ write $\U(q)=(\mathbb Z/q\mathbb Z)^\times$, with $\U(1)$ trivial. All logarithms are natural. An ordered partition of a set $T$ of $m$ distinct primes into $r$ blocks of sizes $k_j$ determines the squarefree factors $F_j$ of $P_T=\prod_{p\in T}p$. A prime may occur in only one block. The partition is uniform among the $\binom{m}{k_1,\ldots,k_r}$ possibilities. Set $H=\prod_{j=1}^r\U(q_j)$, $M=\LCM(q_1,\ldots,q_r)$, and $g=\GCD(q_1,\ldots,q_r)$. A character of $\U(q_j)$ is identified with its pullback $\widetilde\chi_j$ to $G=\U(M)$; the reduction is surjective. The common diagonal subgroup $\Delta_g\subseteq\widehat H$ is the collection of tuples whose pullbacks coincide. It is naturally isomorphic to $\widehat{\U(g)}$: the intersection of the pullback subgroups is the pullback from $g$, by the nested character groups at each prime power.

\begin{theorem}[Growing-label minority-density factorization law]\label{thm:main}
Assume Dirichlet GRH. Fix $0<a\le1$, $A>0$, and put $d=\lfloor1/a\rfloor$. There exist $b,c,C,K>0$ depending only on $a,A$ such that for every sufficiently large $Q$ and integer
\[
2\le R\le c\frac{\log Q}{\log\log Q}
\]
there is a set $S$ of distinct primes in $[K\log^2Q,2K\log^2Q]$ with $\abs S\le CR\log Q$ satisfying the following simultaneously. Choose $2\le r\le R$, arbitrary $q_j\le Q$, set $U=\{p\in S:p\nmid M\}$, and choose any $T\subseteq U$ with $m=\abs T\ge a\abs U$. There is a character subgroup $E\le\widehat H$, independent of quotas, such that
\[
\Delta_g\le E,\qquad [E:\Delta_g]\le d^{r-1}.
\]
Let $\varpi:H\to\widehat E$ be evaluation. For every integer quota vector with $\sum k_j=m$ and $k_j\ge b\log Q$, let $\mu$ be the law of $(F_j\bmod q_j)_j$ and $\nu=\varpi_*\mu$. Then for every $u\in H$,
\begin{align}
\nu(\varpi(u))=0&\quad\Longleftrightarrow\quad \mu(u)=0,\label{eq:supp}\\
\mu(u)&=\frac{\abs E}{\abs H}\nu(\varpi(u))(1+\theta_u),\qquad \abs{\theta_u}\le Q^{-A},\label{eq:law}
\end{align}
where $\theta_u$ can be zero on the empty fibres. Balanced quotas always satisfy the lower bound. At most $d^{r-1}$ quotient values occur, all lying above the compulsory total-product congruence modulo $g$. In particular the fraction of the compatible tuples attained is at least $d^{1-r}$. The bound on the quotient index and this support fraction are sharp (in the worst case over the moduli and deletions).
\end{theorem}

When $a>1/2$, $d=1$ and the only constraint is $\prod_j u_j=P_T\pmod g$. When $a\le1/2$, the finite extra quotient is necessary: the surviving primes may all lie in one coset of a small-index subgroup. Crucially, even as the product group can have order $\exp(\Theta(r\log Q))$, the extra quotient has size $d^{r-1}=Q^{o(1)}$. Formula \eqref{eq:law} counts factorizations after multiplication by the multinomial coefficient. The statement makes no hypothesis of independence of zeros.

\section{A single small-bias pool}
We use the following form of the article's prime-selection lemma (label \texttt{dir:lem:oneset}); a proof is included for clarity. Constants can be enlarged freely, and all onset thresholds below may depend on the fixed parameters.
\begin{lemma}[Uniform small bias]\label{lem:pool}
Assume Dirichlet GRH. Fix $B\ge1$ and $0<\eta_0<1$. There are $K,C_1,c_1>0$ such that, writing $L=\log Q$, for any sufficiently large $Q$ and any real $N_0$ with
\[
 C_1L\le N_0\le c_1 L^2/\log L,
\]
there is a set $S$ of $n$ distinct primes in $[K L^2,2K L^2]$, $N_0\le n\le2N_0$, with
\[
\abs{\sum_{p\in S}\chi(p)}\le\eta_0n
\]
for every nonprincipal character modulo any integer $h\le Q^B$.
\end{lemma}
\begin{proof}
Fix $W\in C_c^\infty((1,2))$ nonnegative with integral one, and write $x=KL^2$ and $w_p=(\log p)W(p/x)$. By the prime number theorem $W_x=\sum_p w_p=x+o(x)$. For a primitive nonprincipal character $\psi$ of conductor $f$, the smoothed explicit formula, using the Mellin transform $\widetilde W(s)=\int_0^\infty W(t)t^{s-1}\,dt$, gives
\[
\sum_{n\ge1}\Lambda(n)\psi(n)W(n/x)
=-\sum_\rho x^\rho\widetilde W(\rho)+O_W(x^{1/4}\log(2f)).
\]
Here zeros are counted with multiplicity. This is obtained by shifting the Mellin integral of $-L'/L$ from $\Re s=2$ to $1/4$ (the line is separated from the nontrivial zeros by GRH); the standard unit-height zero count $O(\log(f(2+|t|)))$ and the corresponding logarithmic derivative bound on that line control the new integral. See the classical explicit formula and zero-count estimates in \cite{mv} and the article's proof of \texttt{dir:lem:oneset}. The rapid decay $\widetilde W(1/2+it)\ll_W(1+|t|)^{-3}$, together with the same zero count, gives $O_W(\sqrt{x}\log(2f))$ for the zero sum. Removing prime powers costs $O_W(\sqrt{x}\log^2(2x))$; passage to a character induced modulo $h$ costs at most $\|W\|_\infty\sum_{p\mid h}\log p\le\|W\|_\infty\log h$ in the prime sum. Consequently for all nonprincipal $\chi$ of modulus $h\le Q^B$,
\[
\frac{\abs{\sum_p w_p\chi(p)}}{W_x}
\ll_W \frac{B\log Q+\log^2(2x)}{\sqrt{x}}+\frac{B\log Q}{x}.
\]
Choose $K$ large depending on $B,\eta_0,W$, then $Q$ large: the right side is at most $\eta_0/4$.

Independently select each prime with probability $\lambda w_p/W_x$, where $\lambda=3N_0/2$. Since $\max w_p\ll_W\log x$ and $W_x\sim x$, these are at most one if $c_1$ is small. The expected sample size is $\lambda$, and the expected complex character sum has modulus at most $\eta_0\lambda/4$. Bernstein's inequality applied separately to the real and imaginary parts, using total variance $\le\lambda$, shows that the probability of an error exceeding $\eta_0\lambda/4$ in a particular complex sum is $\le4e^{-c_2\eta_0^2\lambda}$ for an absolute $c_2>0$. The sample size is in $[N_0,2N_0]$ with failure probability $\le2e^{-c_3\lambda}$. There are at most $\sum_{h\le Q^B}\varphi(h)\le Q^{2B}$ characters. For sufficiently large $C_1$ the union of failure events has probability less than one. On a successful outcome each sum has modulus at most $\eta_0\lambda/2\le\eta_0 n$. This proves existence, with no repeated primes.
\end{proof}

\section{The bounded obstruction group}
The following finite argument is the first modification of the article's minority theorem. We do not need small bias for every character modulo $M$, which could be as large as $Q^R$.
\begin{lemma}[Local words suffice]\label{lem:words}
Let $G$ be a finite abelian group, and let $U$ be a finite labelled list of elements of $G$. Let $T\subseteq U$ have density at least $a$. Let $\mathcal D\subseteq\widehat G$ be closed under inversion, and suppose that distinct characters represented by words of length at most $d=\lfloor1/a\rfloor$ in $\mathcal D$ have inner product of modulus at most $\eta$ in $L^2(U)$, with the uniform probability measure. For sufficiently small $\tau,\eta>0$ depending only on $a$, the subgroup $J$ generated by those $\chi\in\mathcal D$ with $\abs{\E_T\chi}>1-\tau$ has order at most $d$.
\end{lemma}
\begin{proof}
If $t$ distinct words $\chi_i$ have $\abs{\E_T\chi_i}\ge u$, apply the Gram-matrix bound $\| (\langle\chi_i,\chi_j\rangle)_ {i,j}\|\le1+(t-1)\eta$ to the indicator of $T$. With $\alpha=\abs T/\abs U\ge a$ this gives
\[
 t\alpha^2u^2\le\alpha(1+(t-1)\eta).
\]
Every generator with bias $>1-\tau$ differs from a unit complex constant by $L^2(T)$ norm $<\sqrt{2\tau}$. This follows from $\min_{|c|=1}\E|\chi-c|^2=2(1-|\E\chi|)$. Telescoping products, a word of length at most $d$ differs from a constant in norm at most $d\sqrt{2\tau}$ and hence has bias at least $u=1-d^2\tau$. Choose $\tau>0$ so small that $u>0$ and $(d+1)au^2>1$, possible because $(d+1)a>1$. Then take $\eta>0$ with $(d+1)au^2>1+d\eta$. The inequality forbids $d+1$ distinct words in the word ball of radius $d$ generated by the high-bias characters. A word ball of a generating set grows by at least one at each step until it equals the group. Since this ball has size at most $d$, the whole group has size at most $d$.
\end{proof}

Apply this with $G=\U(M)$ and $\mathcal D$ consisting of all ratios $\widetilde\chi_i\widetilde\chi_j^{-1}$ of characters pulled back from any two moduli (including equal indices). For distinct characters represented by words of length $\le d$, their quotient factors through an integer $h\mid M$ of size $h\le Q^{4d}$, and is nonprincipal there. Indeed each factor of a word uses at most two original moduli; the quotient of two words uses at most $4d$ of them. No dependence on $R$ occurs in this conductor bound.

Select $S$ from Lemma~\ref{lem:pool} with $B=4d$, $\eta_0=\eta/4$ and $N_0=C_0RL$, where $C_0$ is a constant to be fixed. At most $v\le r L/\log x$ selected primes divide $M$, because $\prod_{j=1}^r q_j\le Q^r$ and each selected prime is at least $x=KL^2$. Hence $v/n\le1/(C_0\log x)=o(1)$ uniformly. Removing these primes changes the mean of any of the relevant nonprincipal characters to one of modulus at most $(\eta_0 n+v)/(n-v)\le\eta$, since all remaining primes are units. The preceding lemma now produces $J$ of size $D\le d$, depending on $T$ and the moduli but not on quotas. Whenever a ratio in $\mathcal D$ is outside $J$, its bias on $T$ is at most $1-\tau$.

Define
\begin{equation}\label{eq:E}
E=\{(\chi_1,\ldots,\chi_r)\in\widehat H:\ 
\widetilde\chi_j\widetilde\chi_1^{-1}\in J\text{ for every }j\}.
\end{equation}
This is a subgroup; the ratio map into $J^{r-1}$ is a homomorphism with kernel $\Delta_g$. Therefore $[E:\Delta_g]\le D^{r-1}\le d^{r-1}$. Evaluation $\varpi$ is surjective by finite abelian duality. All partition vectors restrict to the same character of $\Delta_g$, determined by $P_T\pmod g$, so at most $[E:\Delta_g]$ fibres can be occupied.

\section{A transportation-cycle compression lemma}
This is the key new finite ingredient. Its conclusion is linear in the number of labels despite the exponentially many possible quotient values. We give an independent abstract formulation, valid for every finite abelian group.
\begin{lemma}[Type-product repair]\label{lem:repair}
Let $B_0$ be a finite abelian group of order $D$ and exponent $e$; use multiplicative notation. Each of $m$ labelled objects has a type $t\in B_0$. Assign them to $r$ labelled blocks of fixed sizes $k_1,\ldots,k_r$. The type signature of an assignment is the $r$-tuple of products of types within its blocks. Given any two assignments $X,Y$ of those sizes, there is an assignment $Z$ with the same signature as $Y$, obtained from $X$ by moving at most
\begin{equation}\label{eq:repair}
H_0=D^2(e-1)r
\end{equation}
objects to different blocks. In particular $H_0\le D^3r$. If $D=1$, no moves are required.
\end{lemma}
\begin{proof}
Let $A=(a_{tj})$ and $B=(b_{tj})$ be the $D\times r$ tables of type counts for $X,Y$. Row sums (type totals) and column sums (block sizes) agree. The integer matrix $B-A$ has every row and column sum zero. Form a directed bipartite graph on the $D$ types and $r$ labels: a positive entry $b_{tj}-a_{tj}$ gives an arc from type $t$ to label $j$ of that multiplicity; a negative entry gives an arc from label $j$ to type $t$ of multiplicity $a_{tj}-b_{tj}$. All vertices are balanced. Decompose this nonnegative integral circulation into simple directed cycles with positive integer coefficients $w_\ell$. To do so, follow positive arcs until a vertex repeats, take the resulting simple cycle, subtract its minimum arc multiplicity, and repeat. Every step erases at least one arc and adds none. There are at most $Dr$ initial arcs, so at most $Dr$ cycles, each of length $\le2D$ and changing at most $D$ object placements per unit coefficient.

Replace each coefficient $w_\ell$ by $w'_\ell=w_\ell\bmod e\in\{0,\ldots,e-1\}$. Add these reduced cycle increments to $A$ to form a table $C$. It has the same margins as $A$ and $B$; each cell lies between the corresponding entries of $A$ and $B$, because every cycle increment on that cell has the sign of $B-A$ and was only reduced. In particular all entries are nonnegative integers. Moreover $C\equiv B\pmod e$ entrywise. Hence the type signature computed from $C$ equals that from $B$, since every type has order dividing $e$.

For each type independently, move objects from surplus cells of $A-C$ to deficit cells; this is possible because row sums agree. The number of moved objects is the sum of the positive parts of $A-C$, at most $D\sum_\ell w'_\ell\le D^2(e-1)r$. All labels retain their quotas. This constructs $Z$ and proves the claim.
\end{proof}

\begin{corollary}[No exponentially thin occupied fibre in the growing regime]\label{cor:mass}
Under the uniform fixed-quota allocation, any positive mass of a type signature, or of any function of that signature, is at least $(1+mr)^{-H_0}$, with $H_0$ as in \eqref{eq:repair}.
\end{corollary}
\begin{proof}
Fix one assignment in the target fibre. Every allocation can be moved into this fibre in Hamming distance at most $H_0$ by Lemma~\ref{lem:repair}. A given allocation has at most $\sum_{j=0}^{\min(H_0,m)}\binom mj(r-1)^j\le(1+mr)^{H_0}$ allocations at that distance (ignoring quotas only enlarges the bound; for $H_0=0$ it is one). Cover the entire allocation space by such balls about the target fibre. Dividing their count by the total number of allocations proves the assertion.
\end{proof}

For our application choose $B_0=\widehat J$ and type $t_p$ the evaluation $\eta\mapsto\eta(p)$ for $\eta\in J$. For an $E$-character tuple $\chi$ let $\eta_j=\widetilde\chi_j\widetilde\chi_1^{-1}\in J$. For the projection of an assignment, evaluation at $\chi$ equals
\[
\prod_j\widetilde\chi_j(F_j)=\widetilde\chi_1(P_T)\prod_{j=2}^r\prod_{p\in T_j}\eta_j(p).
\]
The notation evaluates $F_j$ modulo $M$ before taking its $q_j$ residue. Thus the entire projection is a function of the type signature and the fixed total product. Corollary~\ref{cor:mass} applies to $\nu$. If $r\le L/\log L$, $m\le2C_0RL$, then $\log(1+mr)\le4\log L$ for all sufficiently large $Q$ (since $R\le L/\log L$ and $C_0$ is fixed). Hence every occupied projection satisfies
\begin{equation}\label{eq:mass}
\nu(v)\ge(1+mr)^{-d^3r}\ge Q^{-4d^3}.
\end{equation}
Unlike the cycle argument on the full projection group, this exponent is independent of $r$ in the stated regime. It is immaterial whether every type actually occurs.

\section{Fourier decay outside the obstruction group}
Let $\chi=(\chi_j)\in\widehat H$. Set $z_j(p)=\widetilde\chi_j(p)$ and $w_j=k_j/m$. The characteristic function $\widehat\mu(\chi)=\E\prod_{j}\chi_j(F_j)$ has an elementary torus bound. Assign each prime independently to label $j$ with probability $w_j$, then condition on the counts being $k_j$. The conditioning event has probability at least $(m+1)^{1-r}$. Indeed the count vector $k$ is a mode: from any different count vector $v$, moving one count from $v_i>k_i$ to $v_j<k_j$ multiplies its multinomial probability by $v_i k_j/((v_j+1)k_i)\ge1$. At most $(m+1)^{r-1}$ count vectors are possible. The unconditioned numerator is the coefficient of $\prod_j Z_j^{k_j}$ in
\[
\prod_{p\in T}\left(\sum_j w_j Z_j z_j(p)\right).
\]
On the unit torus the squared modulus of each factor is $1-2\sum_{i<j}w_iw_j(1-\Re(Z_i\overline{Z_j}z_i(p)\overline{z_j(p)}))$. Apply $1-y\le e^{-y}$, take square roots, multiply and use the coefficient bound by torus integration. Discard the nonnegative losses for pairs in the same coset of $J$. Pairs in different cosets have ratio in $\mathcal D\setminus J$, so their loss summed over primes is at least $\tau m$, uniformly in the torus variables. Consequently
\begin{equation}\label{eq:fourier}
|\widehat\mu(\chi)|\le(m+1)^{r-1}\exp\left(-\tau m\sum_{C<C'}w_Cw_{C'}\right),
\end{equation}
where $C$ indexes occupied cosets of $J$ among the pullbacks and $w_C=\sum_{j\in C}w_j$. If $C_0$ has largest total weight and $t$ labels lie outside it, then
\[
\sum_{C<C'}w_Cw_{C'}=\tfrac12(1-\sum_Cw_C^2)\ge\tfrac12(1-w_{C_0})\ge \frac{tk_*}{2m},
\qquad k_*=\min_jk_j.
\]
If $\chi\notin E$, then $t\ge1$. The number of such tuples with a chosen outside-label set of size $t$ is at most $Q\,d^{r-1}Q^t$: choose an anchor character from one of the inside labels (at most $Q$ choices), at most $d$ choices for the pullback of each remaining inside label in its coset, and at most $Q$ choices for each outside label. Pullback is injective for each label. Choosing any largest coset and then its outside-label set is sufficient; overcounting is harmless. Combining with \eqref{eq:fourier}, and setting $z=Qe^{-\tau k_*/2}$, gives
\begin{equation}\label{eq:tail}
\sum_{\chi\notin E}|\widehat\mu(\chi)|\le Qd^{r-1}(m+1)^{r-1}\bigl((1+z)^r-1\bigr).
\end{equation}
For $r\le L/\log L$ and $m\le 2C_0RL$, we have $d^{r-1}\le Q$, $m+1\le L^3$ and $(m+1)^{r-1}\le Q^3$ for large $Q$. Fix
\[
b=\frac{2(A+4d^3+10)}{\tau}.
\]
Then $z\le Q^{1-\tau b/2}$, $rz\le1$ for large $Q$, and $(1+z)^r-1\le e^{rz}-1\le2rz$. Using $r\le Q$, \eqref{eq:tail} yields
\begin{equation}\label{eq:tailsmall}
\sum_{\chi\notin E}|\widehat\mu(\chi)|\le 2Q^{7-\tau b/2}\le Q^{-A-4d^3-2}.
\end{equation}
This is uniform in all quotas, deletions, and moduli. In particular it dominates the potentially tiny positive quotient masses of \eqref{eq:mass} by the requisite relative error.

\section{Proof of the main theorem and consequences}
Choose $\tau,\eta$ as in Lemma~\ref{lem:words} and $b$ as above. Apply Lemma~\ref{lem:pool} with $B=4d$, $\eta_0=\eta/4$ and take $C_0\ge\max(C_1,4(b+1)/a)$. Choose $0<c\le\min(1,c_1/C_0)$, so that $N_0=C_0RL$ meets the lemma's range. Let $C=2C_0$. After exclusion of primes dividing $M$, $m\ge a(n-v)\ge a n/2$ for large $Q$, uniformly, and hence $\lfloor m/r\rfloor\ge bL$. This verifies existence of balanced quotas without making the theorem vacuous.

By Fourier inversion on $H$, the $E$-terms sum exactly to $|E|\nu(\varpi(u))$ when the probability is multiplied by $|H|$. Therefore
\[
\abs{|H|\mu(u)-|E|\nu(\varpi(u))}\le\sum_{\chi\notin E}|\widehat\mu(\chi)|.
\]
If $\nu(\varpi(u))=0$, then $\mu(u)=0$ by nonnegativity and the definition of pushforward. If it is positive, divide by $|E|\nu(\varpi(u))\ge Q^{-4d^3}$, using \eqref{eq:mass} and \eqref{eq:tailsmall}. This proves \eqref{eq:law} with even a $Q^{-2}$ spare factor. Since the resulting multiplicative lower bound is positive, it proves \eqref{eq:supp} too. The order of $E$ and the bound on occupied projections were proved in Section~4. A compatible tuple lies in a fibre of the restriction to $\Delta_g$; the full compatible set has size $|H|/|\Delta_g|$. Every occupied $E$-fibre has size $|H|/|E|$, so the fraction attained is at least $1/[E:\Delta_g]\ge d^{1-r}$.

\paragraph{Sharpness.} For $d\ge2$, choose a fixed prime $\ell\equiv1\pmod d$ (Dirichlet's theorem suffices) and put all $q_j=\ell$. In $U$ retain the most populous coset of the subgroup of index $d$ in $\U(\ell)$. Its density is at least $1/d\ge a$. The $d$ characters of this quotient all have bias one and occur as permitted ratios, so they are contained in $J$; since $|J|\le d$, they equal $J$. All moduli being equal, the ratio map in \eqref{eq:E} is onto $J^{r-1}$, giving equality in the index bound. Each quota fixes the corresponding quotient coset of its factor; thus exactly one $E$-projection occurs. By \eqref{eq:supp} the fraction attained is exactly $d^{1-r}$. For $d=1$, the extra index and fraction are both one. This completes the proof of Theorem~\ref{thm:main}.

\begin{corollary}[Sharp order of the number of factors at this height]\label{cor:order}
At prime height $O(\log^2Q)$, requiring each factor to contain at least $b\log Q$ distinct primes makes $r=O_b(\log Q/\log\log Q)$ necessary. Theorem~\ref{thm:main} attains this order at every fixed positive retained density.
\end{corollary}
\begin{proof}
The number of available primes at this height is $O(\log^2Q/\log\log Q)$ by the prime number theorem. The disjoint blocks consume at least $rb\log Q$ of them. The matching order of growth is furnished by choosing $R=\lfloor c\log Q/\log\log Q\rfloor$ and balanced quotas in the theorem. This is an order-of-growth statement for this quota scale, not an optimal leading constant or a lower bound for all possible weaker notions of coverage.
\end{proof}

\begin{corollary}[Compressed exact support test]\label{cor:algorithm}
For given $S,T,q_j,k_j$ as in the theorem, the occupied fibres and their exact masses can be determined by a dynamic program using at most
\[
D^{r-1}\prod_{j=1}^{r-1}(k_j+1)\le d^{r-1}(m+1)^{r-1}=Q^{O(1)}
\]
 states per layer, where $D=|J|$. This is independent of $|H|$; the full compatible group may be superpolynomial in $Q$.
\end{corollary}
\begin{proof}
Once $J$ is found by enumerating the relevant pair ratios and generating a subgroup, compute the type of each prime. Process the primes in any order, assigning each to any block whose quota is not exhausted. Store the counts and products of types of the first $r-1$ blocks; the last block's count and product are determined by the total count and type product of the processed prefix. Accumulate integer counts in these states, with the initial empty allocation of weight one. After processing all primes, retain states with the desired quotas. The explicit evaluation formula in Section~5 maps each type signature to its $E$-projection; add their counts and divide by the multinomial coefficient. To evaluate a proposed residue tuple, enumerate $E$ by choosing a first character (at most $Q$) and at most $D$ compatible pullbacks for each remaining label, hence at most $Qd^{r-1}$ tuples. Compare its values with the occupied projections. By Theorem~\ref{thm:main} this is exact, even when a compatible tuple is excluded by the minority obstruction. The state bound counts group operations and states, not a fully specified bit-complexity analysis of exact algebraic character evaluation or of constructing the prime pool.
\end{proof}

\paragraph{Verification and limits.} The companion Python file implements the cycle compression over finite abelian groups (using cyclic groups and direct products of cyclic groups for tests), checks the count-table invariants, and compares the exact type-signature dynamic program against exhaustive enumeration on small instances. These are finite consistency checks, not a proof of the analytic estimate. The paper provides the human-readable proof; no Lean formalization of the new result or rerun of the upstream formal development is claimed. The prime pool construction is existential and randomized with positive success probability, not a numerically certified universal pool for a particular $Q$. Constants are not optimized. In particular the theorem does not solve a new nonabelian GRH or claim optimal height for smaller quotas. We have not independently reviewed all 914 results of the merged article; the novelty comparison is against its directly relevant factorization theorems and the targeted published literature.

\begin{thebibliography}{9}\small
\bibitem{vendor} ProveIt, \href{https://github.com/VladimirReshetnikov/ProveIt/tree/main/Analysis/GeneralizedRiemannHypothesis/vendor/grh-lean}{Vendored grh-lean (readable edition)}, README; upstream revision \href{https://github.com/lsernau/grh-lean/tree/72fead102f663d5c7a407964a13dc0cf9d977207}{\texttt{72fead102f6}}.
\bibitem{statements} grh-lean, \href{https://github.com/lsernau/grh-lean/blob/72fead102f663d5c7a407964a13dc0cf9d977207/grh/src/Statements.lean}{\texttt{Statements.lean}}. Vendored blob \texttt{1cd1bce93344065936b2b392b23306254e0785ef}.
\bibitem{main} grh-lean, \href{https://github.com/lsernau/grh-lean/blob/72fead102f663d5c7a407964a13dc0cf9d977207/grh/src/Main.lean}{\texttt{Main.lean}}. Vendored blob \texttt{1adf4187a938510363cdb0dde0a6d628ff39c600}.
\bibitem{ineq} grh-lean, \href{https://github.com/lsernau/grh-lean/blob/72fead102f663d5c7a407964a13dc0cf9d977207/grh/src/New/Reduction/MainInequality.lean}{\texttt{New/Reduction/MainInequality.lean}}.
\bibitem{oneside} grh-lean, \href{https://github.com/lsernau/grh-lean/blob/72fead102f663d5c7a407964a13dc0cf9d977207/grh/src/New/Reduction/OneSidedBound.lean}{\texttt{New/Reduction/OneSidedBound.lean}}.
\bibitem{article} ProveIt contributors, \emph{Beyond the Riemann hypothesis for Dirichlet $L$-functions: consequences, transfers and limits}, 26 November 2026, \href{https://github.com/VladimirReshetnikov/ProveIt/tree/main/Analysis/GeneralizedRiemannHypothesis/Article}{merged article}, especially \texttt{src.tex} and \texttt{intro.tex}.
\bibitem{dir} Same, \href{https://github.com/VladimirReshetnikov/ProveIt/blob/main/Analysis/GeneralizedRiemannHypothesis/Article/sections/dir.tex}{\texttt{sections/dir.tex}}, subsection ``Simultaneous congruences for the factors of a squarefree product'', labels \texttt{dir:thm:factorization}, \texttt{dir:thm:minority}, \texttt{dir:lem:partition}, \texttt{dir:lem:oneset}. Blob \texttt{da58d506e24a05e1b888da2802398f277d9ea727}.
\bibitem{mv} H. L. Montgomery and R. C. Vaughan, \emph{Multiplicative Number Theory I: Classical Theory}, Cambridge University Press, 2007, Ch.~13.
\bibitem{bp} A. R. Booker and C. Pomerance, Squarefree smooth numbers and Euclidean prime generators, \emph{Proc. Amer. Math. Soc.} \textbf{145} (2017), 5035--5042. \href{https://arxiv.org/abs/1607.01557}{arXiv:1607.01557}.
\bibitem{kst} O. Klurman, I. E. Shparlinski, J. Ter\"av\"ainen, On the generation of multiplicative groups by small primes, \href{https://arxiv.org/abs/2609.25393}{arXiv:2609.25393} (2026).
\bibitem{hughes} S. D. Hughes, Few small primes suffice to cover unit groups, \href{https://arxiv.org/abs/2610.04678}{arXiv:2610.04678} (2026).
\end{thebibliography}
\end{document}
